% Part: first-order-logic
% Chapter: completeness
% Section: complete-sets

% Definition of complete consistent sets. Properties of complete sets required
% for completeness proved are in provability.tex in the
% chapter on the proof system used.

\documentclass[../../../include/open-logic-section]{subfiles}

\begin{document}

\iftag{FOL}
      {\olfileid{fol}{com}{ccs}}
      {\olfileid{pl}{com}{ccs}}

\olsection{\usetoken{P}{sentence}ના પૂર્ણ સુસંગત ગણો}

\begin{defn}[પૂર્ણ ગણ]
\ollabel{def:complete-set} !!{sentence}sનો ગણ~$\Gamma$ \emph{!!{complete}}
છે જો અને માત્ર જો કોઈ પણ !!{sentence}~$!A$ માટે કાં તો $!A \in
\Gamma$ અથવા $\lnot !A \in \Gamma$.
\end{defn}

\begin{explain}
વાક્યોના !!^{complete} ગણો એકેય પ્રશ્ન અનુત્તરિત રાખતા નથી. કોઈ
પણ !!{sentence}~$!A$ માટે~$\Gamma$, $!A$ સાચું છે કે ખોટું તે ``કહે''
છે. !!{complete} ગણોનું મહત્વ પૂર્ણતા પ્રમેયની સાબિતી પૂરતું મર્યાદિત
નથી. દાખલા તરીકે, !!{complete} અને સ્વયંસિદ્ધીકરણીય સિદ્ધાંત હંમેશાં
નિર્ણેય હોય છે.
\end{explain}

\begin{explain}
!!^{complete} સુસંગત ગણો પૂર્ણતાની સાબિતીમાં મહત્વના છે, કારણ કે
!!{sentence}sનો દરેક સુસંગત ગણ~$\Gamma$ કોઈ !!a{complete} સુસંગત
ગણ~$\Gamma^*$માં સમાયેલો છે તેની ખાતરી આપી શકાય છે. !!^a{complete}
સુસંગત ગણમાં દરેક !!{sentence}~$!A$ માટે કાં તો~$!A$ અથવા તેનો નિષેધ
$\lnot !A$ હોય છે, પરંતુ બંને નહિ. ખાસ કરીને આ વાત
\iftag{FOL}{પરમાણુક !!{sentence}s}{વિધાનચલો} માટે સાચી છે. તેથી
\iftag{FOL}{!!{constant}s વડે યોગ્ય રીતે વિસ્તૃત ભાષાના}{} કોઈ
!!a{complete} સુસંગત ગણમાંથી એવી
\iftag{FOL}{!!a{structure}}{!!a{valuation}} રચી શકાય છે, જેમાં
\iftag{FOL}{!!{predicate}sનું અર્થઘટન}{વિધાનચલોને નિયુક્ત સત્યમૂલ્ય}
~$\Gamma^*$માં કયાં \iftag{FOL}{પરમાણુક !!{sentence}s}{!!{propositional variable}s}
છે તે પ્રમાણે વ્યાખ્યાયિત થાય છે. પછી આ
\iftag{FOL}{!!{structure}}{!!{valuation}}~$\Gamma^*$નાં બધાં
!!{sentence}sને (અને તેથી~$\Gamma$માંનાં બધાં વાક્યોને પણ) સાચાં બનાવે
છે તેમ બતાવી શકાય છે. આ છેલ્લી હકીકતની સાબિતી માટે $\lnot !A \in
\Gamma^*$ જો અને માત્ર જો $!A \notin \Gamma^*$, $(!A \lor !B) \in
\Gamma^*$ જો અને માત્ર જો $!A \in \Gamma^*$ અથવા $!B \in \Gamma^*$,
વગેરે જરૂરી છે.
\end{explain}

આગળ આપણે ઘણી વાર~$\Proves$ના સ્વવાચકતા, એકદિશવર્ધિતા અને પરંપરિતતાના
ગુણધર્મો અવ્યક્ત રીતે વાપરીશું (જુઓ
\iftag{FOL}{%
  \tagrefs{prfSC/{fol:seq:ptn:sec},prfND/{fol:ntd:ptn:sec},prfAX/{fol:axd:ptn:sec},prfTab/{fol:tab:ptn:sec}}}{%
  \tagrefs{prfSC/{pl:seq:ptn:sec},prfND/{pl:ntd:ptn:sec},prfAX/{pl:axd:ptn:sec},prfTab/{pl:tab:ptn:sec}}}).

\begin{prop}
\ollabel{prop:ccs}
ધારો કે~$\Gamma$ !!{complete} અને સુસંગત છે. તો:
\begin{enumerate}
\item \ollabel{prop:ccs-prov-in} જો $\Gamma \Proves !A$, તો $!A \in
  \Gamma$.

\tagitem{prvAnd}{\ollabel{prop:ccs-and} $!A \land !B \in \Gamma$
  જો અને માત્ર જો $!A \in \Gamma$ અને $!B \in \Gamma$ બંને.}{}

\tagitem{prvOr}{\ollabel{prop:ccs-or} $!A \lor !B \in \Gamma$ જો અને
  માત્ર જો કાં તો $!A \in \Gamma$ અથવા $!B \in \Gamma$.}{}

\tagitem{prvIf}{\ollabel{prop:ccs-if} $!A \lif !B \in \Gamma$ જો અને
  માત્ર જો કાં તો $!A \notin \Gamma$ અથવા $!B \in \Gamma$.}{}
\end{enumerate}
\end{prop}

\begin{proof}
આગળના બધા કિસ્સા માટે ધારો કે~$\Gamma$ !!{complete} અને સુસંગત છે.
\begin{enumerate}
\item જો $\Gamma \Proves !A$, તો $!A \in \Gamma$.

ધારો કે $\Gamma \Proves !A$. પ્રતિપક્ષે ધારો કે $!A \notin \Gamma$.
$\Gamma$ !!{complete} હોવાથી $\lnot !A \in \Gamma$. હવે
\iftag{FOL}{%
  \tagrefs{prfSC/{fol:seq:prv:prop:explicit-inc},
    prfND/{fol:ntd:prv:prop:explicit-inc},
    prfAX/{fol:axd:prv:prop:explicit-inc},
    prfTab/{fol:tab:prv:prop:explicit-inc}}}{%
  \tagrefs{prfSC/{pl:seq:prv:prop:explicit-inc},
    prfND/{pl:ntd:prv:prop:explicit-inc},
    prfAX/{pl:axd:prv:prop:explicit-inc},
    prfTab/{pl:tab:prv:prop:explicit-inc}}}
મુજબ~$\Gamma$ અસુસંગત છે. આ~$\Gamma$ સુસંગત છે એ પૂર્વધારણાનો
વિરોધાભાસ છે. તેથી $!A \notin \Gamma$ હોઈ શકે નહિ, એટલે $!A \in
\Gamma$.

\tagitem{defAnd}{}{%
\iftag{probAnd}{પ્રશ્ન.}{%
$!A \land !B \in \Gamma$ જો અને માત્ર જો $!A \in \Gamma$ અને $!B
\in \Gamma$ બંને:

આગળની દિશા માટે ધારો કે $!A \land !B \in \Gamma$. તો
\iftag{FOL}{%
  \tagrefs{prfSC/{fol:seq:ppr:prop:provability-land},
    prfND/{fol:ntd:ppr:prop:provability-land},
    prfAX/{fol:axd:ppr:prop:provability-land},
    prfTab/{fol:tab:ppr:prop:provability-land}}}{%
  \tagrefs{prfSC/{pl:seq:ppr:prop:provability-land},
    prfND/{pl:ntd:ppr:prop:provability-land},
    prfAX/{pl:axd:ppr:prop:provability-land},
    prfTab/{pl:tab:ppr:prop:provability-land}}}, મુદ્દા~(1) મુજબ
$\Gamma \Proves !A$ અને $\Gamma \Proves !B$. તેથી
\olref{prop:ccs-prov-in} મુજબ $!A \in \Gamma$ અને $!B \in \Gamma$,
જેમ જોઈતું હતું.

ઊલટી દિશા માટે $!A \in \Gamma$ અને $!B \in \Gamma$ લો. હવે
\iftag{FOL}{%
  \tagrefs{prfSC/{fol:seq:ppr:prop:provability-land},
    prfND/{fol:ntd:ppr:prop:provability-land},
    prfAX/{fol:axd:ppr:prop:provability-land},
    prfTab/{fol:tab:ppr:prop:provability-land}}}{%
  \tagrefs{prfSC/{pl:seq:ppr:prop:provability-land},
    prfND/{pl:ntd:ppr:prop:provability-land},
    prfAX/{pl:axd:ppr:prop:provability-land},
    prfTab/{pl:tab:ppr:prop:provability-land}}}, મુદ્દા~(2) મુજબ
$\Gamma \Proves !A \land !B$. તેથી \olref{prop:ccs-prov-in} મુજબ
$!A \land !B \in \Gamma$.}}

\tagitem{defOr}{}{%
\iftag{probOr}{પ્રશ્ન.}{%
પહેલાં બતાવીએ કે જો $!A \lor !B \in \Gamma$, તો કાં તો $!A \in
\Gamma$ અથવા $!B \in \Gamma$. ધારો કે $!A \lor !B \in \Gamma$,
પણ $!A \notin \Gamma$ અને $!B \notin \Gamma$. $\Gamma$ !!{complete}
હોવાથી $\lnot !A \in \Gamma$ અને $\lnot !B \in \Gamma$. હવે
\iftag{FOL}{%
  \tagrefs{prfSC/{fol:seq:ppr:prop:provability-lor},
    prfND/{fol:ntd:ppr:prop:provability-lor},
    prfAX/{fol:axd:ppr:prop:provability-lor},
    prfTab/{fol:tab:ppr:prop:provability-lor}}}{%
  \tagrefs{prfSC/{pl:seq:ppr:prop:provability-lor},
    prfND/{pl:ntd:ppr:prop:provability-lor},
    prfAX/{pl:axd:ppr:prop:provability-lor},
    prfTab/{pl:tab:ppr:prop:provability-lor}}}, મુદ્દા~(1) મુજબ
$\Gamma$ અસુસંગત છે, જે વિરોધાભાસ છે. તેથી કાં તો $!A \in \Gamma$
અથવા $!B \in \Gamma$.

ઊલટી દિશા માટે ધારો કે $!A \in \Gamma$ અથવા $!B \in \Gamma$.
\iftag{FOL}{%
  \tagrefs{prfSC/{fol:seq:ppr:prop:provability-lor},
    prfND/{fol:ntd:ppr:prop:provability-lor},
    prfAX/{fol:axd:ppr:prop:provability-lor},
    prfTab/{fol:tab:ppr:prop:provability-lor}}}{%
  \tagrefs{prfSC/{pl:seq:ppr:prop:provability-lor},
    prfND/{pl:ntd:ppr:prop:provability-lor},
    prfAX/{pl:axd:ppr:prop:provability-lor},
    prfTab/{pl:tab:ppr:prop:provability-lor}}}, મુદ્દા~(2) મુજબ
$\Gamma \Proves !A \lor !B$. તેથી \olref{prop:ccs-prov-in} મુજબ
$!A \lor !B \in \Gamma$, જેમ જોઈતું હતું.}}

\tagitem{defIf}{}{%
\iftag{probIf}{પ્રશ્ન.}{%
આગળની દિશા માટે ધારો કે $!A \lif !B \in \Gamma$, અને પ્રતિપક્ષે ધારો
કે $!A \in \Gamma$ અને $!B \notin \Gamma$. આ પૂર્વધારણાઓ મુજબ
$!A \lif !B \in \Gamma$ અને $!A \in \Gamma$. હવે
\iftag{FOL}{%
  \tagrefs{prfSC/{fol:seq:ppr:prop:provability-lif},
    prfND/{fol:ntd:ppr:prop:provability-lif},
    prfAX/{fol:axd:ppr:prop:provability-lif},
    prfTab/{fol:tab:ppr:prop:provability-lif}}}{%
  \tagrefs{prfSC/{pl:seq:ppr:prop:provability-lif},
    prfND/{pl:ntd:ppr:prop:provability-lif},
    prfAX/{pl:axd:ppr:prop:provability-lif},
    prfTab/{pl:tab:ppr:prop:provability-lif}}}, મુદ્દા~(1) મુજબ
$\Gamma \Proves !B$. પરંતુ પછી \olref{prop:ccs-prov-in} મુજબ
$!B \in \Gamma$, જે $!B \notin \Gamma$ એ પૂર્વધારણાનો વિરોધાભાસ છે.

ઊલટી દિશા માટે પહેલાં $!A \notin \Gamma$ હોય તે કિસ્સો વિચારો.
$\Gamma$ !!{complete} હોવાથી $\lnot !A \in \Gamma$. હવે
\iftag{FOL}{%
  \tagrefs{prfSC/{fol:seq:ppr:prop:provability-lif},
    prfND/{fol:ntd:ppr:prop:provability-lif},
    prfAX/{fol:axd:ppr:prop:provability-lif},
    prfTab/{fol:tab:ppr:prop:provability-lif}}}{%
  \tagrefs{prfSC/{pl:seq:ppr:prop:provability-lif},
    prfND/{pl:ntd:ppr:prop:provability-lif},
    prfAX/{pl:axd:ppr:prop:provability-lif},
    prfTab/{pl:tab:ppr:prop:provability-lif}}}, મુદ્દા~(2) મુજબ
$\Gamma \Proves !A \lif !B$. ફરી \olref{prop:ccs-prov-in} મુજબ
$!A \lif !B \in \Gamma$, જેમ જોઈતું હતું.

હવે $!B \in \Gamma$ હોય તે કિસ્સો વિચારો. ફરી
\iftag{FOL}{%
  \tagrefs{prfSC/{fol:seq:ppr:prop:provability-lif},
    prfND/{fol:ntd:ppr:prop:provability-lif},
    prfTab/{fol:tab:ppr:prop:provability-lif},
    prfAX/{fol:axd:ppr:prop:provability-lif}}}{%
  \tagrefs{prfSC/{pl:seq:ppr:prop:provability-lif},
    prfND/{pl:ntd:ppr:prop:provability-lif},
    prfTab/{pl:tab:ppr:prop:provability-lif},
    prfAX/{pl:axd:ppr:prop:provability-lif}}}, મુદ્દા~(2) મુજબ
$\Gamma \Proves !A \lif !B$. તેથી \olref{prop:ccs-prov-in} મુજબ
$!A \lif !B \in \Gamma$.}}
\end{enumerate}
\end{proof}

\tagprob{FOL}
\begin{prob}
\olref[fol][com][ccs]{prop:ccs}ની સાબિતી પૂરી કરો.
\end{prob}
\tagendprob

\tagprob{notFOL}
\begin{prob}
\olref[pl][com][ccs]{prop:ccs}ની સાબિતી પૂરી કરો.
\end{prob}
\tagendprob

\end{document}
